Pure $O$-sequences, introduced in the context of Stanley’s work on Hilbert functions of monomial level algebras in the late 1970s, have been studied intensively over the last fifteen years \cite{Stanley1978}. They exhibit surprisingly strong positivity phenomena, especially in small codimension and type. Although pure $O$-sequences are not log-concave in general, the case of codimension 3 and type 2 remained open for over a decade \cite{BMMNZ2012,Zanello2024}. The following theorem establishes log-concavity in precisely this remaining case. 

\noindent \textbf{Theorem.}
Every pure $O$-sequence of codimension $3$ and type $2$ is
log-concave.


\noindent\textbf{Terminology.}
A monomial $x_0^{m(0)}x_1^{m(1)}x_2^{m(2)}$ in three variables is
identified with its exponent vector $m=(m(0),m(1),m(2))\in\mathbb N^3$. 
We write $m'\le m$ for coordinate-wise inequality and
$|m|=m(0)+m(1)+m(2)$ for total degree.  A finite order ideal
$\Gamma\subset\mathbb N^3$ is a finite set closed downward under this
order: if $n\in\Gamma$ and $m\le n$, then $m\in\Gamma$.

A monomial $m\in\Gamma$ is maximal if every $n\in\Gamma$ with
$m\le n$ also satisfies $n\le m$.  The ideal is pure if all maximal
monomials have the same total degree; this common degree is the socle
degree.  The type is the number of maximal monomials.  The pure
$O$-sequence of $\Gamma$ is
\[
  h_\Gamma(d)=\#\{m\in\Gamma:|m|=d\}.
\]
For a monomial $g$, we simply write $h_g(d)=\#\{u\le g:|u|=d\}$. The $O$-sequence is log-concave if
\[
  h_\Gamma(i-1)h_\Gamma(i+1)\le h_\Gamma(i)^2
  \qquad(0<i<e),
\]
where $e$ is the socle degree of $\Gamma$.  We write
$\mathbf 1_P$ for the indicator of the condition $P$.

\noindent \textbf{Proof.}
Let $\Gamma$ be a pure order ideal in three variables and suppose that
its two maximal monomials are $m_1\ne m_2$.  Since $\Gamma$ is pure, $|m_1|=|m_2|=e$.
Every element of a finite order ideal lies below a maximal element, so
\[
  \Gamma=\{u:u\le m_1\} \cup \{u:u\le m_2\},
\]
After possibly interchanging $m_1$ and $m_2$, there is a coordinate
$x\in\{0,1,2\}$ such that
\[
  m_2(x)>m_1(x),
  \qquad
  m_2(y)\le m_1(y)\quad(y\ne x).                         \tag{1}
\]
Indeed, the three integers $m_1(i)-m_2(i)$ sum to zero and are not all
zero, hence one sign occurs in a minority position; if necessary we swap
the two monomials so that the minority sign is negative.

Let $c=m_1(x)+1$, and define
\[
  M_{\rm red}(y)=
  \begin{cases}
    m_2(x)-m_1(x)-1,&y=x,\\
    m_2(y),&y\ne x.
  \end{cases}
\]
Then $|M_{\rm red}|=e-c$. For $u\le m_2$, condition (1) implies
\[
  u\not\le m_1
  \quad\Longleftrightarrow\quad
  u(x)\ge c.
\]
If $d\ge c$, the map $u\mapsto u-ce_x$ bijects
$\{u\le m_2:|u|=d,\ u(x)\ge c\}$
with
$\{v\le M_{\rm red}:|v|=d-c\}$.
If $d<c$, the first set is empty.  Hence the $O$ sequence of
$\Gamma$ can be written as 
\[
  H(d):=h_\Gamma(d)
  =
  h_{m_1}(d)+\mathbf 1_{d\ge c}\,h_{M_{\rm red}}(d-c).       \tag{2}
\]
We need to prove that $H(d-1)H(d+1)\le H(d)^2$ holds for $0<d<e$.

For an integer-valued sequence $F:\mathbb N\to\mathbb Z$, define the first and scond difference sequences as
\[
  \Delta F(0)=F(0),
  \qquad
  \Delta F(d)=F(d)-F(d-1)\quad(d>0),
\]
\[
  \Delta^2F(0)=\Delta F(0),
  \qquad
  \Delta^2F(d)=\Delta F(d)-\Delta F(d-1)\quad(d>0).
\]

We first prove an elementary divisor-count formula for a single
monomial. Let $m=(\alpha,\beta,\gamma)$ be a monomial. 
For $p,q\in\mathbb N$, set
\[
  L_{p,q}(k)=\#\{r\in\mathbb N:r<p,\ r\le k,\ k-r<q\}.
\]
Thus $L_{p,q}(k)$ counts the number of pairs $(r,s)$ with
$0\le r<p$, $0\le s<q$, and $r+s=k$.  A direct box count gives
\[
  h_m(t)
  =
  \sum_{k=0}^{t}\mathbf 1_{t-k<\gamma+1}\,
  L_{\alpha+1,\beta+1}(k).
\]
With the convention $\Delta F(0)=F(0)$, the first-difference identities are
\[
  \Delta L_{p,q}(t)
  =
  \mathbf 1_{t<p}
  +
  \mathbf 1_{t<q}
  -
  \mathbf 1_{t<p+q}
  \qquad(t\ge0),
\]
and hence
\[
  \Delta h_m(t)
  =
  L_{\alpha+1,\beta+1}(t)
  -
  \mathbf 1_{t\ge \gamma+1}\,
  L_{\alpha+1,\beta+1}(t-\gamma-1)
  \qquad(t\ge 0).
\]
Taking one more difference gives, for every $t>0$,
\begin{align}
  \Delta^2 h_m(t)
  &=
  \mathbf 1_{t\le\alpha}
  +
  \mathbf 1_{t\le\beta}
  +
  \mathbf 1_{t\le\gamma}
  -
  \mathbf 1_{t\le\alpha+\beta+1}
  -
  \mathbf 1_{t\le\beta+\gamma+1}       \notag\\
  &\hspace{2.5cm}
  -
  \mathbf 1_{t\le\gamma+\alpha+1}
  +
  \mathbf 1_{t\le\alpha+\beta+\gamma+2}.              \tag{3}
\end{align}
In particular, $\Delta^2h_m(t)\le 1$. 
The same bound holds for the shifted sequence
$t\longmapsto \mathbf 1_{t\ge c}h_{M_{\rm red}}(t-c),
$
because below the shift its second difference is zero, and from the shift
onward it is the second difference of the principal sequence of
$M_{\rm red}$.  Therefore $\Delta^2H(t)\le 2$.                                      

Formula (3) also implies the following useful dichotomy.  If
$t>0$ and $\Delta^2h_m(t)=1$, then either
\[
  t\le m(0),\qquad t\le m(1),\qquad t\le m(2),              \tag{L}
\]
or
\[
  t\ge m(0)+m(1)+2,\qquad
  t\ge m(1)+m(2)+2,\qquad
  t\ge m(0)+m(2)+2,\qquad
  t\le |m|+2.                                             \tag{U}
\]
We refer to these as the lower and upper alternatives. For $t=0$, the lower alternative is automatic. 

We shall also use the following elementary consequences of the same box
count. 
\begin{enumerate}
    \item If $t\le m(i)$ for all $i$, then
\[
  2h_m(t)=(t+1)(t+2),
  \qquad
  \Delta h_m(t)=t+1.                                      \tag{5}
\]
\item If $t\ge m(i)+m(j)$ for every pair $\{i,j\}$,
and $t\le |m|+2$,
then,
\[
  2h_m(t)=(|m|+1-t)(|m|+2-t).                              \tag{6}
\]
\item If $t\ge m(i)+m(j)+1$ for every pair $\{i,j\}$,
and $t\le |m|+2$, then 
\[
  \Delta h_m(t)=-(|m|+2-t).                                \tag{7}
\]
\item For all $t$,
\[
  2h_m(t)\le (|m|-t+1)(|m|-t+2).                           \tag{8}
\]
\item If $2t\ge |m|+1$, then
\[
  \Delta h_m(t)\le0.                                      \tag{9}
\]
\end{enumerate}
Now fix $0<d<e=|m_1|=|m_2|$. Using $H(d)=H(d-1)+\Delta H(d)$
and $\Delta H(d+1)=\Delta H(d)+\Delta^2H(d+1)$
gives
\[
  H(d)^2-H(d-1)H(d+1)
  =
  (\Delta H(d))^2-H(d-1)\Delta^2H(d+1).                   \tag{10}
\]
Since $H(d-1)\ge0$ and $\Delta^2H(d+1)\le2$, it is enough to prove the following two statements:
\[
  \Delta^2H(d+1)=1
  \quad\Longrightarrow\quad
  H(d-1)\le(\Delta H(d))^2,                               \tag{11}
\]
and
\[
  \Delta^2H(d+1)=2
  \quad\Longrightarrow\quad
  2H(d-1)\le(\Delta H(d))^2.                              \tag{12}
\]
Indeed, if $\Delta^2H(d+1)\le0$, equation (10) is immediate; if it is
$1$ or $2$, then (11) or (12) applies.

Write $f(t)=h_{m_1}(t)$ and $g(t)=\mathbf 1_{t\ge c}h_{M_{\rm red}}(t-c)$, so that $H=f+g$.
Since $\Delta^2f\le1$ and $\Delta^2g\le1$, if
$\Delta^2H(d+1)=1$, then either
\[
  \Delta^2f(d+1)=1,
  \qquad
  \Delta^2g(d+1)=0,                                      \tag{A}
\]
or
\[
  \Delta^2f(d+1)=0,
  \qquad
  \Delta^2g(d+1)=1.                                      \tag{B}
\]
If $\Delta^2H(d+1)=2$, then
\[
  \Delta^2f(d+1)=\Delta^2g(d+1)=1.
  \tag{C}
\]

\medskip
\noindent\textbf{Case A1: (A)+(L) for $f$.}
Suppose equation (A) holds and the lower alternative (L) holds for $f$ with $t=d+1$. Then in particular $d+1\le m_1(x)<c$, thus $g(d-1)=g(d)=0$, hence $\Delta g(d)=0$ and therefore
\[
  H(d-1)=f(d-1),
  \qquad
  \Delta H(d)=\Delta f(d).
\]
By (5), $2f(d-1)=d(d+1)$ and $\Delta f(d)=d+1$. 
Therefore
\[
  H(d-1)=f(d-1)\le(d+1)^2=(\Delta f(d))^2=(\Delta H(d))^2.
\]
This proves (11) in this case.

\medskip
\noindent\textbf{Case A2: (A)+(U) for $f$}. Assume equation (A) holds and the upper alternative (U) holds for $f$ with $t=d+1$. Let us introduce the shorthand notation $\rho=e-d$.
Then (U) gives $\rho\ge -1$.  Applying
(6) and (7) gives
\[
  2f(d-1)=(\rho+2)(\rho+3),
  \qquad
  \Delta f(d)=-(\rho+2).                                    \tag{13}
\]
The shifted summand satisfies
\[
  \Delta g(d)\le0,
  \qquad
  2g(d-1)\le (\rho+2)(\rho+3),
  \qquad
  2g(d)\le (\rho+1)(\rho+2).        \tag{14}
\]
Indeed, the two inequalities for $g(d-1)$ and
$g(d)$ follow from the universal upper bound (8),
applied to the shifted sequence.  For the sign of $\Delta g(d)$, if
$d<c$, then $\Delta g(d)=0$.  If $d\ge c$, then the upper alternative
for $f$ implies
\[
  d\ge m_1(x)+m_1(y)+1,
  \qquad
  d\ge m_1(x)+m_1(z)+1,
\]
where $\{x,y,z\}=\{0,1,2\}$.  Adding these inequalities gives
\[
  2d\ge 2m_1(x)+m_1(y)+m_1(z)+2=e+m_1(x)+2.
\]
Since $c=m_1(x)+1$ and $|M_{\rm red}|=e-c$, this is equivalent to
$2(d-c)\ge |M_{\rm red}|+1$.  Thus the midpoint monotonicity (9) applies
to the shifted $M_{\rm red}$-sequence and gives $\Delta g(d)\le0$.

To prove (11), by (13) it remains to show
\[
  f(d-1)+g(d-1)\le\bigl(-(\rho+2)+\Delta g(d)\bigr)^2.
\]
By (13) and (14), this follows from
\[
  (\rho+2)(\rho+3)+2g(d-1)
  \le
  2\bigl(-(\rho+2)+\Delta g(d)\bigr)^2.                    \tag{15}
\]
Since $g(d)=g(d-1)+\Delta g(d)$, the third inequality in (14) gives
\[
  2g(d-1)
  \le
  (\rho+1)(\rho+2)-2\Delta g(d).
\]
Therefore the difference between the right-hand side and the left-hand
side of (15) is at least
\begin{align*}
&2\bigl(-(\rho+2)+\Delta g(d)\bigr)^2
  -(\rho+2)(\rho+3)
  -\bigl((\rho+1)(\rho+2)-2\Delta g(d)\bigr)  \\
&\qquad
=
2\Delta g(d)\bigl(\Delta g(d)-2(\rho+2)+1\bigr).
\end{align*}
This is nonnegative because $\Delta g(d)\le0$, $\rho\ge-1$, and hence
\[
  \Delta g(d)-2(\rho+2)+1\le0.
\]
Thus (11) holds in this case, and we completed the (A) case.

For the two cases when (B) holds, that is $\Delta^2g(d+1)=1$, we use the following
shift reduction.  If
\[
  \Delta^2H(d+1)=1
  \qquad\text{and}\qquad
  H(d+1)\le(\Delta H(d+1))^2,
\]
then
\[
  H(d-1)\le(\Delta H(d))^2.                                \tag{16}
\]
Indeed, $\Delta H(d)=\Delta H(d+1)-1$ and $H(d-1)=H(d+1)-\Delta H(d+1)-\Delta H(d)$.
Therefore
\[
  H(d-1)
  =
  H(d+1)-2\Delta H(d+1)+1
  \le
  (\Delta H(d+1))^2-2\Delta H(d+1)+1
  =
  (\Delta H(d))^2.
\]

\medskip
\noindent\textbf{Case B1: (B)+(L) for $g$.}
Assume equation (B) holds and the lower alternative (L) holds for $g$ with $t=d+1$. Let $D=d+1$. Since $\Delta^2g(D)=1$, necessarily $D\ge c$, and $\Delta^2h_{M_{\rm red}}(D-c)=1$.
In the lower subcase $D-c\le M_{\rm red}(i)$ for every $i$. Then $D-c+1\ge 1$, and the lower formula (5) gives
\[
  2g(D)=(D-c+1)(D-c+2),
  \qquad
  \Delta g(D)=D-c+1.                                           \tag{17}
\]
We next prove two bounds for $f$ at $D$:
\[
  f(D)\le (D+1)\Delta f(D),
  \qquad
  \Delta f(D)=c.                                       \tag{18}
\]
First, since $\Delta^2f(D)=0$, the explicit formula (3), together with
the lower inequalities for $M_{\rm red}$, implies
\[
  D\le m_1(0)+m_1(1)+1,
  \qquad
  D\le m_1(1)+m_1(2)+1,
  \qquad
  D\le m_1(0)+m_1(2)+1.                                  \tag{19}
\]
Indeed, the two pair inequalities involving $x$ follow directly from
$D-c\le M_{\rm red}(y)\le m_1(y)$ for $y\ne x$.  The remaining pair
inequality follows from (3): if it failed, then the corresponding two
single indicators would be zero; since $D\ge c=m_1(x)+1$, the
$x$-single indicator is also zero, and the right-hand side of (3)
would be negative, contradicting $\Delta^2f(D)=0$.

Thus, in (3) for $f$ at $D$, all three pair indicators and the triple
indicator are equal to $1$.  Since $\Delta^2f(D)=0$, exactly two of
the three single indicators are equal to $1$.  Moreover the
$x$-single indicator is zero, because $D\ge m_1(x)+1$.  Hence the
two remaining single indicators are equal to $1$.

It follows that for every $1\le k\le D$, the same two single
indicators are still equal to $1$, while the three pair indicators and
the triple indicator are also equal to $1$.  Therefore
\[
  \Delta^2f(k)\ge 2-3+1=0
  \qquad(1\le k\le D).
\]
Thus the first differences $\Delta f(k)$ are nondecreasing on
$\{0,1,\ldots,D\}$.  Since
\[
  f(D)=\sum_{j=0}^{D}\Delta f(j),
\]
we get
\[
  f(D)\le (D+1)\Delta f(D).
\]
This proves the first inequality in (18).

For the second inequality in (18), relabel the two coordinates different
from $x$ as $y,z$, and write
\[
  A=m_1(y)+1,
  \qquad
  B=m_1(z)+1,
\]
and recall $c=m_1(x)+1$. Using the box-count formula with $x$ as the third coordinate, the
first-difference formula gives
\[
  \Delta f(D)=L_{A,B}(D)-L_{A,B}(D-c),
\]
because $D\ge c$.  From the preceding paragraph,
$D<A$ and $D<B$, hence
\[
  L_{A,B}(D)=D+1,
  \qquad
  L_{A,B}(D-c)=D-c+1.
\]
Therefore
\[
  \Delta f(D)=(D+1)-(D-c+1)=c,
\]
proving (18). Using (17) and (18), we have
\[
  2H(D)
  =
  2f(D)+2g(D)
  \le
  2\Delta f(D)(D+1)+(D-c+1)(D-c+2).
\]
Also, by (17),
\[
  \Delta H(D)=\Delta f(D)+\Delta g(D)
  =
  \Delta f(D)+(D-c+1).
\]
Thus it is enough to show
\[
  2\Delta f(D)(D+1)+(D-c+1)(D-c+2)
  \le
  2\bigl(\Delta f(D)+D-c+1\bigr)^2.
\]
Since $D+1=(D-c+1)+c$, the difference between the right-hand side and
the left-hand side is
\[
  2\Delta f(D)\bigl(\Delta f(D)-c\bigr)
  +2\Delta f(D)(D-c+1)
  +(D-c+1)(D-c).
\]
This is nonnegative because $\Delta f(D)\ge c$ and $D-c+1\ge1$.
Therefore $H(D)\le(\Delta H(D))^2$.
By the shift reduction (16), this proves (11).

\medskip
\noindent\textbf{Case B2: (B)+(U) for $g$.}
Again let $D=d+1$ and put $\rho=e-D$.  The upper alternative for
$h_{M_{\rm red}}$ at $D-c$ gives $\rho\ge -2$.  Since
$|M_{\rm red}|=e-c$, formulas (6) and (7) give
\[
  2g(D)=(\rho+1)(\rho+2),
  \qquad
  \Delta g(D)=-(\rho+2).                                  \tag{20}
\]
For the first summand, the universal upper bound (8) gives
\[
  2f(D)\le(e-D+1)(e-D+2)=(\rho+1)(\rho+2).                 \tag{21}
\]

It remains to justify the sign of $\Delta f(D)$.  Since $g$ is in the
upper alternative at $D$, we have
\[
  D-c\ge M_{\rm red}(0)+M_{\rm red}(1)+2,\ 
  D-c\ge M_{\rm red}(1)+M_{\rm red}(2)+2,\ D-c\ge M_{\rm red}(0)+M_{\rm red}(2)+2.
\]
Adding these three inequalities gives
\[
  3(D-c)\ge 2|M_{\rm red}|+6=2(e-c)+6.
\]
Hence 
\[
  2(D-c)
  \ge
  \frac{2}{3}\bigl(2|M_{\rm red}|+6\bigr)
  =
  |M_{\rm red}|+1+\frac{|M_{\rm red}|+9}{3}
  >
  |M_{\rm red}|+1,
\]
because $|M_{\rm red}|\ge0$.
Therefore
\[
  2D=2(D-c)+2c\ge (e-c+1)+2c=e+c+1\ge e+1.
\]
By midpoint monotonicity (9), applied to $f=h_{m_1}$, we get $\Delta f(D)\le0$. Using (20), (21), we get 
\[
  2H(D)\le2(\rho+1)(\rho+2),
  \qquad
  \Delta H(D)=\Delta f(D)-(\rho+2).
\]
Since $\rho\ge-2$ and $\Delta f(D)\le0$, we have
\[
  (\rho+1)(\rho+2)\le\bigl(-\Delta f(D)+\rho+2\bigr)^2.
\]
giving the desired $H(D)\le(\Delta H(D))^2$. By the shift reduction (16), this proves (11).

\medskip
\noindent\textbf{Case C: (C) holds.}
So here $\Delta^2H(d+1)=2$ and hence $\Delta^2f(d+1)=\Delta^2g(d+1)=1$.
Since $\Delta^2g(d+1)=1$, necessarily $d+1\ge c$.
Therefore the lower alternative (L) for $f$ is impossible, because it
would imply $d+1\le m_1(x)<c$.
Hence $f$ is in its upper alternative, that is, (U) holds for $f$. The lower alternative (L) for $g$ is also impossible. Indeed, if $d+1-c\le M_{\rm red}(i)$ for every $i$,
then, choosing any $y\ne x$, we get
\[
  d+1
  \le
  m_1(x)+1+M_{\rm red}(y)
  =
  m_1(x)+1+m_2(y)
  \le
  m_1(x)+1+m_1(y).
\]
But the upper alternative for $f$ gives $d+1\ge m_1(x)+m_1(y)+2$, a contradiction.  Thus $g$ is also in its upper alternative, and (U) holds for $g$.
Put $\rho=e-d$ again. Since $|M_{\rm red}|=e-c$, we also have $|M_{\rm red}|-(d-c)=e-d=\rho$.
The upper formulas (6) and (7), applied to $f$ and to the shifted
summand $g$, give
\[
  2f(d-1)=(\rho+2)(\rho+3),
  \qquad
  \Delta f(d)=-(\rho+2),
\]
and
\[
  2g(d-1)=(\rho+2)(\rho+3),
  \qquad
  \Delta g(d)=-(\rho+2).
\]
Therefore
\[
  2H(d-1)=2(\rho+2)(\rho+3),
  \qquad
  \Delta H(d)=-2(\rho+2).
\]
Since the upper alternative for $f$ gives $\rho\ge-1$, we have
\[
  (\Delta H(d))^2-2H(d-1)
  =
  4(\rho+2)^2-2(\rho+2)(\rho+3)
  =
  2(\rho+2)(\rho+1)
  \ge0.
\]
which proves (12). This completes the proof of (11) and (12), and hence by (10) $H(d-1)H(d+1)\le H(d)^2$.
Since $0<d<e$ was arbitrary and $H$ is the pure $O$-sequence of
$\Gamma$ by (2), the theorem follows. \hfill $\Box$

